%%=============================================================================
%% Inleiding
%%=============================================================================

\chapter{\IfLanguageName{dutch}{Inleiding}{Introduction}}%
\label{ch:inleiding}

%De inleiding moet de lezer net genoeg informatie verschaffen om het onderwerp te begrijpen en in te zien waarom de onderzoeksvraag de moeite waard is om te onderzoeken. In de inleiding ga je literatuurverwijzingen beperken, zodat de tekst vlot leesbaar blijft. Je kan de inleiding verder onderverdelen in secties als dit de tekst verduidelijkt. Zaken die aan bod kunnen komen in de inleiding~\autocite{Pollefliet2011}:

%\begin{itemize}
  %\item context, achtergrond
  %\item afbakenen van het onderwerp
  %\item verantwoording van het onderwerp, methodologie
  %\item probleemstelling
  %\item onderzoeksdoelstelling
  %\item onderzoeksvraag
  %\item \ldots
%\end{itemize}

\section{\IfLanguageName{dutch}{Probleemstelling}{Problem Statement}}%
\label{sec:probleemstelling}

%Uit je probleemstelling moet duidelijk zijn dat je onderzoek een meerwaarde heeft voor een concrete doelgroep. De doelgroep moet goed gedefinieerd en afgelijnd zijn. Doelgroepen als ``bedrijven,'' ``KMO's'', systeembeheerders, enz.~zijn nog te vaag. Als je een lijstje kan maken van de personen/organisaties die een meerwaarde zullen vinden in deze bachelorproef (dit is eigenlijk je steekproefkader), dan is dat een indicatie dat de doelgroep goed gedefinieerd is. Dit kan een enkel bedrijf zijn of zelfs één persoon (je co-promotor/opdrachtgever).%/

Graafmachines, hijskranen en betonpompen zijn cruciaal voor de dagelijkse activiteiten van kleine bouwbedrijven en hebben een directe impact op zowel de productiviteit als de kosten van
een bouwproject waaraan het bedrijf werkt \autocite{Mikulic2024}. Als deze machines plotseling defect raken, resulteert dit vaak in stilstand, vertragingen en aanzienlijke kosten voor reparaties van de machines \autocite{ErRatby2025}. In tal van kleine bouwondernemingen wordt onderhoud tegenwoordig nog altijd voornamelijk preventief uitgevoerd op vaste tijdstippen. Volgens
het onderzoek van \textcite{Morganti2024} en \textcite{Deepak2025} tonen zij aan dat deze methode niet altijd effectief is en dit kan resulteren in onnodige onderhoudskosten of onvoorziene stilstanden van de machines.

\section{\IfLanguageName{dutch}{Onderzoeksvraag}{Research question}}%
\label{sec:onderzoeksvraag}

%Wees zo concreet mogelijk bij het formuleren van je onderzoeksvraag. Een onderzoeksvraag is trouwens iets waar nog niemand op dit moment een antwoord heeft (voor zover je kan nagaan). Het opzoeken van bestaande informatie (bv. ``welke tools bestaan er voor deze toepassing?'') is dus geen onderzoeksvraag. Je kan de onderzoeksvraag verder specifiëren in deelvragen. Bv.~als je onderzoek gaat over performantiemetingen, dan 

De belangrijkste onderzoeksvraag is: Op welke manier kan een machine learningmodel, gebaseeerd op sensorgegevens, worden gebruikt om storingen in bouwmachines bij een kleine bouwonderneming nauwkeurig te voorspellen en het onderhoud effectiever te organiseren?
Eerst wordt onderzocht welke soorten sensordata en onderhoudsproblemen van belang zijn voor bouwmachines \autocite{Mikulic2024, Nayak2024}. Welke technische beperkingen er zijn binnen kleine bouwondernemingen \autocite{Elahi2023}. 
Zoals \textcite{Carvalho2019}, \textcite{Roehrich2024}, \textcite{Li2025} en \textcite{Wang2025} ook in hun onderzoek beschrijven, kan daarna pas onderzocht worden welke machine learning algoritmes geschikt zijn voor voorspellend onderhoud en op welke manier deze correct kunnen worden beoordeeld.

\section{\IfLanguageName{dutch}{Onderzoeksdoelstelling}{Research objective}}%
\label{sec:onderzoeksdoelstelling}

%Wat is het beoogde resultaat van je bachelorproef? Wat zijn de criteria voor succes? Beschrijf die zo concreet mogelijk. Gaat het bv.\ om een proof-of-concept, een prototype, een verslag met aanbevelingen, een vergelijkende studie, enz.

Deze bachelorproef heeft als doel een soort van proof-of-concept voorspellend onderhoudsmodel te creëren en te evalueren voor de gekozen casus. Dit wordt ook aangevuld met een praktisch
adviesrapport om het model te integreren en uit te voeren binnen het bedrijf. De bachelorproef is pas succesvol, zoals \textcite{Raffaele2024} en \textcite{Mahale2025} ook beschrijven in hun onderzoek, als het model technisch helemaal goed functioneert en een duidelijke waarde toevoegt aan het onderhoudsbeleid van het bouwbedrijf.

\section{\IfLanguageName{dutch}{Opzet van deze bachelorproef}{Structure of this bachelor thesis}}%
\label{sec:opzet-bachelorproef}

% Het is gebruikelijk aan het einde van de inleiding een overzicht te
% geven van de opbouw van de rest van de tekst. Deze sectie bevat al een aanzet
% die je kan aanvullen/aanpassen in functie van je eigen tekst.

De rest van deze bachelorproef is als volgt opgebouwd:

In Hoofdstuk~\ref{ch:stand-van-zaken} wordt een overzicht gegeven van de stand van zaken binnen het onderzoeksdomein, op basis van een literatuurstudie.

In Hoofdstuk~\ref{ch:methodologie} wordt de methodologie toegelicht en worden de gebruikte onderzoekstechnieken besproken om een antwoord te kunnen formuleren op de onderzoeksvragen.

% TODO: Vul hier aan voor je eigen hoofstukken, één of twee zinnen per hoofdstuk

In Hoofdstuk~\ref{ch:conclusie}, tenslotte, wordt de conclusie gegeven en een antwoord geformuleerd op de onderzoeksvragen. Daarbij wordt ook een aanzet gegeven voor toekomstig onderzoek binnen dit domein.